\clearpage
\section{Radioactive Decay}
\subsection{5 minutes}
The decay rate is known so the different minutes are independent meaning
\begin{equation}
    N_5=5N_1
\end{equation}
but since \(N_1\) is distribution poisson,
\begin{equation}
    N_1=\lambda\pm\sqrt{\lambda}
\end{equation}
so 
\begin{equation}
    N_5 = 5\lambda\pm\sqrt{5\lambda}=500\pm10\sqrt{5}
\end{equation}
\subsection{Unknown material}
Since we assume the decay is poisson this tells us that 
\begin{equation}
    \lambda=100
\end{equation}
which implies 
\begin{equation}
    \sigma = \sqrt{\lambda}=10
\end{equation}
which is a variance whose source is the measurement. We add this to the variance from the last part which is inherent to the decay and get 
\begin{equation}
    N = 500\pm10\sqrt{5}\pm50
\end{equation}
\subsection{composite material}
Since we measured the number of decays of both materials together and they are independent we measured 
\begin{equation}
    N_i\sim Pois\ins{\lambda_1+\lambda_2}
\end{equation}
so this is not enough to measure each separately. Any attempt to create an equation that separates them will fail, and we'll always have dependence on \(\lambda_1+\lambda_2\).